\documentclass{iopjournal}
\usepackage{subfig}
\usepackage{graphicx}
\usepackage{eurosym}
\usepackage{amsmath}
\usepackage{amssymb}
\usepackage{siunitx}
\usepackage{multirow}

% Custom units not predefined by siunitx
\DeclareSIUnit{\MBq}{MBq}
\DeclareSIUnit{\kBq}{kBq}
\DeclareSIUnit{\MeV}{MeV}
\DeclareSIUnit{\keV}{keV}
\DeclareSIUnit{\Gy}{Gy}
\DeclareSIUnit{\kcps}{kcps}

\graphicspath{{figs/}{figs/crysp/}}

\begin{document}

\articletype{Paper}

\title{An in-room PET for dose verification in proton therapy
based on cryogenic CsI: the CRYSP approach}

\author{J. J. G\'omez-Cadenas$^{1}$ and the CRYSP collaboration$^{1}$}

\affil{$^1$Donostia International Physics Center, San Sebasti\'an / Donostia, E-20018, Spain}

\email{jjgomezcadenas@dipc.org}

\keywords{Proton therapy, range verification, in-room PET, Total Body PET, cryogenic CsI}

\begin{abstract}
 Proton therapy delivers highly conformal dose but is
very sensitive to range uncertainties, motivating in vivo, in situ
verification of the delivered treatment. PET imaging of the
$\beta^{+}$-emitters produced by the beam is the only practical approach,
and among the three operational modalities the in-room configuration offers
the best compromise. We examine the requirements that an in-room PET must
satisfy to be useful as a dosimetric tool and show that an approach based in Cryogenic Cristals for PET (CRYSP) may provide near-to-ideal solutions to this problem. %\textit{Approach.} Starting from published quantitative estimates of the
%induced activity, its temporal evolution and biological washout, we derive
%the detector requirements and compare scintillator technologies for this
%specific, photon-starved measurement. \textit{Main results.} The dominant
%requirement is sensitivity --- driven primarily by the axial field of view
%(AFOV) --- combined with the ability to recover oblique coincidences without
%loss, which in turn demands excellent energy resolution (to reject
%patient scatter) and depth-of-interaction (DOI) capability (to suppress
%parallax). Conventional LYSO scanners are poorly matched: their intrinsic
%$^{176}$Lu radioactivity contaminates a sub-MBq signal, their main asset
%(time-of-flight) shows diminishing returns for this task, and their cost is
%prohibitive for a large AFOV. BGO removes the radioactivity but its poor
%energy resolution and low light yield (precluding monolithic DOI detectors)
%trade against the wrong figures of merit. \textit{Significance.} CRYSP
%combines $<7\%$ energy resolution, millimetre 3D resolution with DOI, and a
%crystal cost up to an order of magnitude below LYSO, while its only weakness
%--- pile-up at high rate --- is irrelevant in the low-activity proton regime.
%It is therefore a strong candidate for an affordable, high-performance
%in-room PET for proton therapy verification.
\end{abstract}

\section{Introduction}
\label{sec:intro}

Proton therapy is among the most precise modalities of external radiation
therapy. Unlike a photon beam, which deposits a high entrance dose that
decreases gradually through the body, a proton beam deposits most of its
energy near the end of its range, at the Bragg peak, and delivers essentially
no dose beyond it~\cite{paganetti2012range}. This depth--dose characteristic
enables highly conformal fields, a much lower integral dose (of order
\SI{60}{\percent} lower than photon therapy), and is particularly valuable for
tumours near critical structures and for paediatric patients, where the risk
of radiation-induced secondary tumours is a concern.

The very feature that makes proton therapy attractive --- the sharp distal
dose fall-off --- also makes it acutely sensitive to range uncertainties.
A range error can leave part of the tumour unirradiated (\emph{under-shooting})
or deliver full dose to healthy tissue distal to the target
(\emph{over-shooting})~\cite{paganetti2012range}. Uncertainties arise in
treatment planning (CT noise and artefacts, partial-volume effects, and
especially the Hounsfield-unit to stopping-power conversion) and during
delivery (setup and positioning errors, organ motion, anatomical changes).
Accurate knowledge of the delivered dose, and ideally a fast feedback signal,
is therefore essential to exploit the full physical advantage of protons and
to set safe planning margins.

Because the proton beam stops completely inside the patient, there is no exit
detector and direct in vivo monitoring is difficult~\cite{zhu2013proton}.
Of the proposed non-invasive techniques, MRI detects physiological changes
that develop over weeks and is not an in situ method, while prompt-gamma
imaging must contend with emissions in the \SIrange{4}{10}{\MeV} range for
which no practical detectors yet exist. The only practical approach for
in vivo, in situ verification is PET imaging of the positron emitters induced
along the beam path by target-fragmentation nuclear
reactions~\cite{zhu2013proton}. The relevant species in soft tissue are
$^{15}$O, $^{11}$C and $^{13}$N (see section~\ref{sec:requirements}). Crucially,
dose deposition (ionisation) and positron-emitter production (nuclear
reactions) are distinct processes: the activity distribution does
not coincide with the dose, so the PET image must be compared against a
predicted activity distribution (Monte Carlo or analytical) rather than
directly with the dose~\cite{zhu2013proton}.

%\begin{figure}[htbp]
%\centering
%\includegraphics[width=\linewidth]{fig4_three_modalities.png}
%\caption{The three operational modalities for PET verification of proton
%therapy (after Zhu \& El~Fakhri~\cite{zhu2013proton}): (a) in-beam PET,
%with detector panels integrated in the beam delivery system; (b) off-line PET,
%where the patient is moved to a nearby PET facility; (c) in-room PET, a
%stand-alone, full-ring scanner in the treatment room that images the patient
%on the treatment couch shortly after irradiation.}
%\label{fig:modalities}
%\end{figure}

Three operational modalities are possible.
% have been pursued, illustrated in
%fig.~\ref{fig:modalities}~\cite{zhu2013proton}.
\emph{In-beam} PET integrates detectors into the beam delivery system,
acquiring during and immediately after irradiation; it offers the best timing
but is expensive and geometrically constrained, typically forcing dual-head
configurations with incomplete angular coverage. \emph{Off-line} PET uses an
established scanner nearby; it is inexpensive but the \SIrange{15}{30}{\minute}
transport delay means the short-lived $^{15}$O has decayed and only $^{11}$C is
measured, with severe washout. \emph{In-room} PET --- a stand-alone,
full-ring scanner positioned in the treatment room --- is the best compromise:
the patient is moved to the scanner without repositioning and imaged a few
minutes after irradiation, retaining abundant $^{15}$O signal while allowing a
full-ring geometry. In the pioneering MGH studies the mean treatment-to-scan
delay was \SI{2.5}{\minute}~\cite{zhu2011monitoring,min2013clinical}, and
in-room PET is identified as the most cost-effective option for most
hospital-based proton centres~\cite{zhu2013proton}.

Using PET as a dosimetric tool in the in-room configuration imposes a specific
set of requirements. The induced activity is of order \SI{1}{\MBq} at its
peak, and about 100-300 kB by the time the PET scan starts. This is roughly three orders of magnitude below a diagnostic PET study, which
translates into only $\sim 10^{5}$--$10^{6}$ collected coincidences for a
typical field~\cite{parodi2008comparison}. This, in turn, requires a scanner with high sensitivity and the capability to use every LOR for image reconstruction. %The system must therefore be
%extremely sensitive; it must acquire quickly (a delay $\lesssim$\,\SI{2}{\minute})
%to capture $^{15}$O; it needs reliable attenuation correction and
%co-registration with the planning CT, which calls for an integrated CT
%component or fiducials; and it benefits from time-resolved acquisition to
%disentangle the isotope species and model washout. 
 
 In this work, we argue that a  large-AFOV scanner with
excellent energy resolution and DOI capability provides an optimal solution to this problem. In particular, we propose an apparatus
based on cryogenic monolithic  crystals --- the CRYSP design~\cite{Soleti:2024flw}. 

The paper is organised as follows. Section~\ref{sec:requirements} reviews the
quantitative features of the measurement --- isotope activities, decay
kinetics and biological washout --- and derives the detector requirements,
showing why conventional LYSO (and, with caveats, BGO) scanners are poorly
suited. Section~\ref{sec:donostia} specialises these to the concrete
installation at Donostia, fixing the modality and the default operating point.
Section~\ref{sec:crysp} summarises the CRYSP design and presents it as a
solution, and section~\ref{sec:extensions} discusses three extensions of the
concept.

\section{The measurement and its requirements}
\label{sec:requirements}

\subsection{Induced activity and its temporal evolution}

During irradiation, target-fragmentation reactions produce positron emitters
along the beam path. For protons only target fragmentation is possible (no
projectile fragmentation), so the relevant channels in soft tissue produce
$^{15}$O ($T_{1/2}=\SI{2.0}{\minute}$), $^{11}$C ($T_{1/2}=\SI{20.4}{\minute}$)
and $^{13}$N ($T_{1/2}=\SI{10.0}{\minute}$), with minor contributions from
$^{30}$P and $^{38}$K~\cite{zhu2013proton}. Because the PET signal is a
mixture of species with very different half-lives, the measurement depends
strongly on the time course of acquisition.

These same reactions carry a biological significance of their own: the
short-range, high-LET secondary particles released in target fragmentation are
believed to be one of the contributions that raise the proton relative
biological effectiveness (RBE) above the constant value of $1.1$ assumed in
current treatment planning~\cite{paganetti2014relative}. Since the primary
protons are only weakly ionising along the entrance and plateau region, the
fragmentation contribution is relatively most relevant there, before the Bragg
peak. RBE variability in proton therapy is an actively debated topic, and
mapping the fragmentation-induced activity probes precisely the reactions
implicated in it --- an additional motivation for in vivo verification beyond
the geometric check of the range.

\begin{figure}[htbp]
\centering
\includegraphics[width=0.8\linewidth]{fig3_relative_contributions.png}
\caption{Activity of the main radionuclide species as a function of time after
irradiation, due to radioactive decay, for a clinical case (production rates
from Parodi et al.~\cite{parodi2008comparison}, as compiled by Zhu \&
El~Fakhri~\cite{zhu2013proton}). The $y$-axis is the activity in MBq and the
$x$-axis the time in minutes.}
\label{fig:activity}
\end{figure}

Figure~\ref{fig:activity} shows the activity, by species, as a function of time
after irradiation. Table~\ref{tab:activity} reads off approximate values at
representative acquisition times. The total peak activity is of order
\SIrange{0.5}{0.6}{\MBq} --- sub-MBq, i.e. hundreds of kBq. The $^{15}$O
component starts at $\sim$\SI{0.5}{\MBq} but decays with a \SI{2}{\minute}
half-life: by \SI{10}{\minute} it is comparable to $^{11}$C, and by
$\sim$\SI{15}{\minute} it has effectively vanished. The $^{11}$C component is
nearly flat ($T_{1/2}=\SI{20}{\minute}$) at
$\sim$\SIrange{0.04}{0.05}{\MBq} and comes to dominate after
$\sim$\SIrange{10}{13}{\minute}. This is precisely why the acquisition instant
matters: an in-room scan at $\sim$\SI{2}{\minute} captures the full
$\sim$\SI{0.35}{\MBq} with $^{15}$O still dominant, whereas an off-line scan at
$\geq$\SI{15}{\minute} sees only $\sim$\SI{0.04}{\MBq}, essentially $^{11}$C
alone --- roughly an order of magnitude less signal.

\begin{table}[htbp]
\centering
\caption{Approximate activity (in MBq) read from fig.~\ref{fig:activity}, by
species and acquisition time. $^{15}$O dominates the first minutes; $^{11}$C
dominates after $\sim$\SIrange{10}{13}{\minute}.}
\label{tab:activity}
\begin{tabular}{lcccc}
\hline
Time & $^{15}$O & $^{11}$C & $^{13}$N & \textbf{Total} \\
     & ($T_{1/2}\!=\!2.0$\,min) & ($20.4$\,min) & ($10.0$\,min) & \\
\hline
$t\approx0$ (end of beam)      & $\sim$0.50 & $\sim$0.05  & $\sim$0.01  & $\sim$\textbf{0.58} \\
$\sim$2 min (in-room)          & $\sim$0.30 & $\sim$0.045 & $\sim$0.01  & $\sim$\textbf{0.35} \\
$\sim$10 min                   & $\sim$0.05 & $\sim$0.04  & $\sim$0.005 & $\sim$\textbf{0.10} \\
$\sim$15--20 min (off-line)    & $\sim$0    & $\sim$0.03  & $\sim$0     & $\sim$\textbf{0.04} \\
$\sim$30 min                   & $\sim$0    & $\sim$0.02  & $\sim$0     & $\sim$\textbf{0.02} \\
$\sim$45 min                   & $\sim$0    & $\sim$0.01  & $\sim$0     & $\sim$\textbf{0.01} \\
\hline
\end{tabular}
\end{table}

\subsection{Decay kinetics and the measured counts}
\label{sec:kinetics}

Each species $j$ decays with constant $\lambda_j=\ln 2/T_{1/2,j}$ (assumed
$100\%$ $\beta^{+}$). For a delivery of duration $t_{\mathrm{irr}}$ at an
approximately constant production rate, followed by a delay $t_{\mathrm{del}}$
before an acquisition window $t_{\mathrm{meas}}$, the number of decays collected
inside the window follows the model of Parodi et al.~\cite{parodi2008comparison}.
For a continuous (cyclotron-like) delivery --- the relevant case for our
installation, see section~\ref{sec:donostia} --- it reduces to a transparent
three-factor expression,
\begin{equation}
N_{j,\mathrm{meas}} =
\underbrace{\frac{P_j}{\lambda_j t_{\mathrm{irr}}}\left(1-e^{-\lambda_j t_{\mathrm{irr}}}\right)}_{\text{nuclei alive at end of beam}}
\;\underbrace{e^{-\lambda_j t_{\mathrm{del}}}}_{\text{decay in transport}}
\;\underbrace{\left(1-e^{-\lambda_j t_{\mathrm{meas}}}\right)}_{\text{fraction decaying in window}},
\label{eq:nmeas}
\end{equation}
where $P_j$ is the integral production of species $j$, proportional to the
delivered dose. The first factor accounts for build-up and decay during
delivery, the second for decay during patient transport, and the third for the
fraction of the surviving nuclei that decay while the scanner is acquiring.
Equation~\eqref{eq:nmeas} makes explicit why a short delay is decisive: for
$^{15}$O ($T_{1/2}=\SI{122}{\second}$) the transport factor
$e^{-\lambda_j t_{\mathrm{del}}}$ has fallen to $\sim$0.5 already at
$t_{\mathrm{del}}=\SI{2}{\minute}$ and to $\sim$0.06 at \SI{15}{\minute}, whereas
for $^{11}$C ($T_{1/2}=\SI{1223}{\second}$) it is still $\sim$0.93 at
\SI{2}{\minute}. The short in-room delay is precisely what preserves the
$^{15}$O signal.

\subsection{Washout and time-resolved acquisition}

Biological washout begins as soon as irradiation starts and progressively
alters both the magnitude and the spatial distribution of the
activity~\cite{parodi2008comparison}. Parodi et al. estimate that washout
reduces the measurable signal by $\sim$\SIrange{30}{40}{\percent} (head) and
$\sim$\SIrange{20}{30}{\percent} (paraspinal) even for in-beam imaging, and by
up to $\sim$\SI{55}{\percent} off-line. The standard three-component model
(fast/medium/slow) is derived from carbon-beam animal studies, assigns the same
parameters to $^{11}$C and $^{15}$O in a given tissue, and ignores local
perfusion; accurate washout parameters are notoriously hard to
obtain~\cite{zhu2013proton}.

A proposed alternative is kinetic modelling of $^{15}$O~\cite{zhu2013proton},
which avoids tabulated washout parameters. Instead of a single static image, the
signal is reconstructed in several time frames (a dynamic acquisition), yielding
a time--activity curve in each voxel. The curve is fitted to a two-parameter rate
equation,
\begin{equation}
\frac{dN}{dt} = R_{\mathrm{prod}} - \left(\lambda + k_{\mathrm{bio}}\right) N,
\label{eq:kinetic}
\end{equation}
where $R_{\mathrm{prod}}$ is the $^{15}$O production rate, which tracks the
delivered dose and range, and the elimination rate is the sum of the known
radioactive constant $\lambda$ and the unknown biological washout
$k_{\mathrm{bio}}$. The fit recovers both rates at once, so the washout is
estimated from the data rather than assumed from a table.

This approach requires $^{15}$O to be abundant enough to build a time--activity
curve. Owing to its \SI{2}{\minute} half-life, $^{15}$O survives only at short
delays: it is present in the in-beam and in-room modalities, but has decayed in
the off-line case, where only $^{11}$C remains. Disentangling the species and
modelling the washout thus requires time-resolved (list-mode) acquisition, with
the counting statistics per frame as the limiting factor.

This is one of the reasons the in-room modality is of interest here. Of the two
configurations in which $^{15}$O is available, in-beam is expensive and
geometrically constrained, while in-room combines a short delay --- so that
$^{15}$O is still abundant --- with a full-ring geometry and a lower cost. In-room
thus appears as a good compromise for preserving the $^{15}$O signal while
providing adequate sensitivity and image quality.

\subsection{Counting statistics: the central constraint}

The decisive figure for detector design is the available number of decays.
For a single \SI{1}{\Gy} field, Parodi et al. find integral yields of order
$10^{7}$--$10^{8}$ decays~\cite{parodi2008comparison}. Assuming \SI{80}{\percent}
attenuation and a \SI{4}{\percent} detection efficiency, this corresponds to
only \numrange{80000}{800000} coincidences over the whole acquisition --- far
below the typical values of nuclear-medicine tracer imaging. As the authors
note, ``even a few percent reduction of the available signal\dots may have
significant consequences on the quality of the reconstructed
images''~\cite{parodi2008comparison}. The measurement is fundamentally
photon-starved.

\subsection{Detector requirements}
\label{sec:reqs}

Because only $10^{5}$--$10^{6}$ coincidences are available, the design is
governed by the need to collect and use as many of the gamma pairs as possible.
This translates into two coupled requirements.

\paragraph{Maximum sensitivity, driven by the AFOV.}
With only $10^{5}$--$10^{6}$ coincidences available, the geometric collection
efficiency is paramount. The dominant lever is the axial field of view, which
sets the solid angle subtended by the activated volume. For a localised source
centred in a scanner of radius $R$ and half-length $L$, the fraction of
back-to-back pairs whose line of response intercepts the detector scales
approximately as $L/\sqrt{L^{2}+R^{2}}$. With $R\simeq\SI{38.7}{\cm}$, this
fraction rises from $\sim$\SI{55}{\percent} at $L=\SI{25.6}{\cm}$
(\SI{0.5}{\metre} AFOV) to $\sim$\SI{80}{\percent} at $L=\SI{51.2}{\cm}$
(\SI{1}{\metre} AFOV): halving the length would cost $\sim$\SI{45}{\percent} of
the relative sensitivity, a large loss in this regime. A large AFOV is
therefore a primary requirement (fig.~\ref{fig:acceptance}).

\begin{figure}[htbp]
\centering
\includegraphics[width=0.65\linewidth]{acceptance.pdf}
\caption{Geometrical acceptance for point-like and extended sources as a
function of the axial length of the scanner, computed at the transverse
centre of a \SI{80}{\cm}-diameter PET (from the CRYSP study~\cite{Soleti:2024flw}).
Most of the gain is realised within the first \SI{100}{\cm} of axial coverage.}
\label{fig:acceptance}
\end{figure}

\paragraph{Control of multiple scattering, to keep oblique coincidences.}
A long AFOV only helps if the oblique coincidences it adds are actually usable.
Two effects degrade them. First, gammas emitted at small angles to the axis
traverse long paths through the patient and are more likely to Compton-scatter
in the body, producing a wrong line of response; this background is suppressed
by selecting a narrow window around \SI{511}{\keV}, whose efficiency improves
with the energy resolution of the detector. Second, gammas entering the crystals
obliquely suffer parallax error unless the depth of interaction is measured. Both
effects are largest at the large axial angles where a long AFOV contributes most.
A detector for this task should therefore combine good energy resolution with DOI
capability, so that the high-angle coincidences are retained. The two
requirements are complementary: the AFOV sets how many pairs are collected, the
energy resolution and DOI how many of them are usable.

\subsection{Conventional scintillators}

Conventional total-body PET scanners are built either with a short AFOV, which
limits the solid angle available for this localised source, or with a large AFOV
at a cost dominated by the scintillator. Two material-specific facts are relevant
to the comparison developed in this work.

LYSO, the standard choice in modern scanners, contains $^{176}$Lu, a naturally
radioactive isotope ($\beta^{-}$ decay with a $\gamma$ cascade) that gives an
intrinsic activity of a few hundred Bq per cm$^{3}$ of
crystal~\cite{yamamoto2005investigation,conti2017physics}. This self-activity is
negligible against the hundreds of MBq of a diagnostic study, but for a sub-MBq,
localised signal it contributes a floor of singles and random coincidences. LYSO
is also a rare earth whose cost dominates that of large-AFOV systems.

BGO has a higher density and effective atomic number, giving a larger
photoelectric fraction and a higher intrinsic sensitivity per unit crystal
volume, and it contains no intrinsic radioactivity. Its light yield, however, is
low (about 9\,000\,photons/MeV), which limits the achievable energy resolution
and, in a monolithic geometry, the DOI reconstruction.

How these characteristics weigh against sensitivity, energy resolution, DOI and
cost for the in-room measurement is what the comparison in this work is intended
to quantify. In the following we consider a scintillator with a high light
yield, good energy resolution and no intrinsic radioactivity: cryogenic CsI.

\section{Installation in Donostia}
\label{sec:donostia}

The proton-therapy facility in Donostia/San Sebasti\'an is built around an IBA
Proteus ONE single-room system. Its accelerator is an S2C2 superconducting
synchrocyclotron: a pulsed machine ($\sim$\SI{1}{\kilo\hertz} pulse
repetition, $\sim$\SI{10}{\micro\second} pulses) that delivers a fixed
\SI{230}{\MeV} beam, degraded to a clinical range of \SIrange{70}{230}{\MeV} by
a variable-energy degrader at the cyclotron exit, with pencil-beam scanning at a
conventional dose rate of $\sim$\SI{2}{\Gy\per\minute}. This concrete setting fixes
both the modality and the operating point of the verification PET.

\subsection{Choice of modality}

Two features of the installation favour the in-room configuration. First, the
compact single-room geometry leaves little space for the dual-head panels of an
in-beam system. Such a configuration, restricted to two opposed heads, covers a
much smaller solid angle than the full ring available to an in-room scanner, with
correspondingly lower sensitivity and limited-angle reconstruction artefacts ---
a drawback for a measurement in which sensitivity is the leading requirement.
Second, the pulsed time structure produces prompt-$\gamma$ and neutron
backgrounds during the beam pulses, which complicate acquisition during
irradiation. Off-line imaging at a remote scanner, on the other hand, incurs the
\SIrange{15}{30}{\minute} delay that suppresses the $^{15}$O signal through the
transport factor of Eq.~\eqref{eq:nmeas}. An in-room scanner --- a stand-alone
full-ring system beside the nozzle, imaging the patient on the couch a short
delay after irradiation --- therefore appears as a good compromise for this
installation.

Because acquisition is entirely post-beam, the pulsed time structure of the S2C2
does not affect the count budget: once the beam is off, the induced activity
decays according to its own physics regardless of how it was delivered. For the
count budget the machine behaves as the continuous (cyclotron) limit, and the
spill/pause bookkeeping needed for in-beam synchrotron acquisition reduces to the
three-factor expression of Eq.~\eqref{eq:nmeas}.

\subsection{Operating point and default parameters}

We adopt the following defaults for the Donostia case, used throughout the
detector study and in the companion simulation chain:
\begin{itemize}
  \item delivered dose $D=\SI{2}{\Gy}$ per field (which sets the integral
  production $P_j$);
  \item irradiation time $t_{\mathrm{irr}}=\SI{2}{\minute}$ ($\approx\SI{120}{\second}$);
  \item treatment-to-scan transport delay $t_{\mathrm{del}}=\SI{2}{\minute}$
  ($\approx\SI{120}{\second}$), consistent with the $\approx\SI{2.5}{\minute}$
  mean of the MGH in-room studies~\cite{zhu2011monitoring};
  \item acquisition window $t_{\mathrm{meas}}=\SI{20}{\minute}$
  ($\SI{1200}{\second}$), as in the clinical in-room example of
  Zhu~\&~El~Fakhri~\cite{zhu2013proton}.
\end{itemize}

Evaluating Eq.~\eqref{eq:nmeas} with these values gives the surviving fractions
of Table~\ref{tab:donostia}. At a delay of \SI{2}{\minute} the $^{15}$O activity
exceeds that of $^{11}$C by a factor of $\approx 4$, so the measurement is
$^{15}$O-dominated. With a peak total activity at the sub-MBq level
(section~\ref{sec:requirements}) and a \SI{2}{\Gy} field, the available
$10^{7}$--$10^{8}$ decays translate, after attenuation and detection efficiency,
into only $\sim 10^{5}$--$10^{6}$ collected coincidences. This photon-starved,
$^{15}$O-rich operating point is what drives the detector requirements of
section~\ref{sec:reqs}, and is the reference scenario reproduced in the companion
simulation study.

\begin{table}[htbp]
\centering
\caption{Surviving fractions per species from Eq.~\eqref{eq:nmeas} at the
Donostia default operating point ($t_{\mathrm{irr}}=t_{\mathrm{del}}=\SI{120}{\second}$,
$t_{\mathrm{meas}}=\SI{1200}{\second}$). The last column is the activity at the
start of acquisition relative to $^{11}$C, assuming comparable production $P_j$;
the actual $P_j$ are provided by the Monte Carlo. The measurement is
$^{15}$O-dominated.}
\label{tab:donostia}
\begin{tabular}{lccccc}
\hline
Species & $T_{1/2}$ (s) & build-up & transport & window & rel.\ initial \\
        &               & $\frac{1-e^{-\lambda t_{\mathrm{irr}}}}{\lambda t_{\mathrm{irr}}}$ & $e^{-\lambda t_{\mathrm{del}}}$ & $1-e^{-\lambda t_{\mathrm{meas}}}$ & activity \\
\hline
$^{15}$O & 122  & 0.73 & 0.51 & 1.00 & \textbf{4.1} \\
$^{13}$N & 598  & 0.93 & 0.87 & 0.75 & 1.8 \\
$^{11}$C & 1223 & 0.97 & 0.93 & 0.49 & 1.0 \\
\hline
\end{tabular}
\end{table}

\section{CRYSP as a solution}
\label{sec:crysp}

CRYSP is a Total-Body PET concept based on large monolithic crystals of
pure caesium iodide (CsI) operated at cryogenic temperature
($\sim$\SIrange{77}{100}{\kelvin})~\cite{Soleti:2024flw}. Cooling pure CsI
shifts its emission towards the near-UV and raises its light yield by a factor
of $\sim$20 to $\sim 10^{5}$ photons/MeV~\cite{mikhailik2015luminescence},
about three times that of LYSO. This high light output underlies the
detector characteristics relevant to the present problem.

\begin{figure}[htbp]
\centering
\includegraphics[width=0.9\linewidth]{cryspDesign.pdf}
\caption{Conceptual design of CRYSP. The monolithic crystals (purple) are bathed
in liquid nitrogen and read out by matrices of SiPMs (green). Thin vacuum
chambers provide thermal insulation; the inner wall facing the patient is at
room temperature, so the bore is a conventional warm tube~\cite{Soleti:2024flw}.}
\label{fig:crysp}
\end{figure}

The baseline scanner, CRYSP1M, has a ring diameter of \SI{77.4}{\cm} and an
AFOV of \SI{102.4}{\cm} (20 rings of 48 crystals each, 960 monolithic crystals
of $48\times48\times37$\,mm$^{3}$, i.e. two radiation lengths). Each crystal is
read out by an $8\times8$ array of SiPMs, for a total of $\sim$61\,440 channels
--- a factor of four fewer than a comparable pixelated system. Its relevant
performance figures, taken from the CRYSP study~\cite{Soleti:2024flw}, map
directly onto the requirements derived above.

\paragraph{Energy resolution.}
Cryogenic CsI achieves an energy resolution of \SI{6.3}{\percent} FWHM at
\SI{511}{\keV}~\cite{Soleti:2024flw,moszynski2005energy}, better than LYSO
($\sim$\SI{10}{\percent}) or BGO. This allows a tight \SIrange{465}{555}{\keV}
window that suppresses patient scatter, and is the property that makes the
high-angle coincidences from a long AFOV usable. Moreover, the excellent energy
resolution enables the recovery of events that Compton-scatter once and are then
photo-absorbed within the same crystal --- up to $\sim$\SI{60}{\percent} of all
events, a larger fraction than in higher-$Z$ materials --- so that fewer pairs
are discarded.

\begin{figure}[htbp]
\centering
\includegraphics[width=0.75\linewidth]{event_display.pdf}
\caption{Charge collected by the SiPM matrix for two interactions at different
depths in a monolithic crystal. The difference in the light pattern is exploited
by a convolutional neural network to reconstruct the interaction vertex,
including the depth of interaction~\cite{Soleti:2024flw}.}
\label{fig:doi}
\end{figure}

\paragraph{Monolithic geometry, DOI and spatial resolution.}
Rather than pixels, CRYSP reads out monolithic crystals and reconstructs the 3D
interaction vertex with a convolutional neural network from the SiPM light
pattern (fig.~\ref{fig:doi}). This yields an effective resolution of
$\sim$\SI{1.7}{\mm} FWHM in all three coordinates, including the depth of
interaction, so that parallax is strongly reduced and the resolution stays
roughly constant with the transaxial angle~\cite{Soleti:2024flw}. This is the
second ingredient needed to keep oblique coincidences, and it meets the
millimetre spatial resolution required for clinically useful range verification.

\paragraph{Sensitivity and cost.}
With its \SI{1}{\metre} AFOV, CRYSP1M provides the large solid angle that the
measurement demands. Its NEMA sensitivity ($\sim$\SI{120}{\kcps\per\MBq}) is somewhat
below LYSO total-body scanners, but the better energy resolution yields
reconstructed images of comparable quality at equal numbers of events, and is
favourable for sources deep inside a scattering medium ---
the situation of an activated volume inside a patient~\cite{Soleti:2024flw}.
Crucially, the cost of monolithic CsI can be up to an order of magnitude below
pixelated LYSO, while cryogenic operation (a simple liquid-nitrogen bath, with
none of the complexity of superconducting magnets) adds less than
\SI{5}{\percent} to the system cost and is transparent to clinical
users~\cite{Soleti:2024flw}. A large, hence sensitive, scanner is therefore
affordable.

\paragraph{No intrinsic radioactivity, and pile-up made irrelevant.}
CsI contains no naturally radioactive isotope, removing the $^{176}$Lu
background floor that handicaps LYSO in a sub-MBq measurement. The principal
weakness of CRYSP --- pile-up arising from the $\sim$\SI{1}{\micro\second} decay
time at high counting rates --- is essentially eliminated in this application:
the proton-induced activity is sub-MBq, orders of magnitude below the rates
at which pile-up becomes significant. CRYSP does not provide TOF; for this
low-rate, scatter-dominated measurement it relies on energy resolution and DOI
rather than on timing, and the relative value of TOF here is one of the questions
the comparison addresses.

\paragraph{Summary.}
CRYSP combines a large AFOV for sensitivity, good energy resolution and DOI to
retain the oblique coincidences that the AFOV provides, millimetre spatial
resolution for range verification, freedom from intrinsic radioactivity, and a
cost compatible with broad clinical deployment, with pile-up --- its main
limitation --- not significant at the low activity of this signal. The
integration of a CT component for attenuation correction and co-registration with
the planning CT is, as in any PET/CT, an engineering addition that does not
involve the cryogenic detector. These properties make CRYSP a candidate worth
evaluating for an affordable in-room PET for proton therapy verification.

\section{CRYSP extensions}
\label{sec:extensions}

The baseline CRYSP detector --- monolithic pure CsI at cryogenic temperature
--- is one point in a family of designs that share the same architecture
(monolithic crystals, SiPM readout, CNN-based 3D vertex reconstruction with DOI,
and freedom from intrinsic radioactivity) but make different trade-offs between
cost, sensitivity and energy resolution. Because the in-room proton measurement
is intrinsically low-rate, the slow scintillation timing common to all these
materials --- and the consequent absence of TOF --- is not a penalty, which
opens the door to options that would be uncompetitive in a high-rate diagnostic
setting. We outline three extensions, spanning an entry-level to a high-end
configuration.

\subsection{Room-temperature CsI(Tl): the entry-level workhorse}
\label{sec:csitl}

The simplest extension removes the cryostat altogether. Thallium-doped CsI,
CsI(Tl), operated at room temperature, has a light yield of
$\sim$54\,000\,photons/MeV --- higher than LYSO and roughly half that of
cryogenic pure CsI --- with the same host crystal and therefore the same
density, effective atomic number and radiation length as the baseline. This
light output is more than sufficient to read out large monolithic crystals and
to reconstruct the 3D interaction vertex, including DOI, with the same neural-network
approach. At this photon level the energy resolution is limited not by
photostatistics but by the intrinsic non-proportionality of the scintillator (the
Fano-like floor), so CsI(Tl) reaches an energy resolution close to that of
cryogenic pure CsI, preserving the scatter rejection the measurement requires.

Its scintillation is slow ($\sim$1\,$\mu$s decay), even slower than cryogenic
pure CsI, but in the low-rate proton regime this carries no penalty: there is no
pile-up and no useful TOF to forgo. The decisive advantage is the removal of
cryogenics: complexity, cost and operational overhead drop substantially while
every other element of the design --- monolithic geometry, DOI, energy
resolution, absence of intrinsic radioactivity --- is retained. The
hygroscopicity of CsI(Tl) is a matter of crystal packaging rather than a
performance limitation, and afterglow is not relevant at these activities. This
makes the uncooled CsI(Tl) scanner an attractive entry-level option: the cheapest
and simplest route to a first deployment.

\subsection{Cryogenic BGO: trading energy resolution for density}
\label{sec:cryobgo}

At the opposite end of the sensitivity argument, cooling BGO to $\sim$\SI{77}{\kelvin}
raises its light yield from $\sim$9\,000\,photons/MeV at room temperature to
$\sim$20\,000\,photons/MeV --- enough, now, to support monolithic crystals with
DOI reconstruction and to improve the energy resolution relative to
room-temperature BGO. The motivation is BGO's high density and effective atomic
number, which give a larger photoelectric fraction and hence a higher intrinsic
sensitivity per unit crystal volume than CsI --- a direct gain on the leading
requirement of the measurement.

The low light yield of BGO noted in section~\ref{sec:requirements} refers to
room-temperature operation; cooling supplies additional light and makes a
monolithic, DOI-capable BGO detector feasible. The timing degrades further at low
temperature, where the slow scintillation component dominates, but this is not
significant in the low-rate regime. The energy resolution of cryogenic BGO
remains below that of pure CsI, so this option favours density-driven sensitivity
over scatter rejection. Like CsI, BGO has no intrinsic radioactivity.

\subsection{A hybrid cryogenic scanner: BGO core, CsI wings}
\label{sec:hybrid}

The two preceding extensions favour opposite properties --- density and
sensitivity (BGO) against energy resolution and DOI (CsI). A third configuration
combines them using the geometry of the measurement: an axially graded scanner
with a BGO core and pure-CsI wings, both cryogenic.

The rationale follows from where each property matters. For a localised source
centred in the AFOV, the photons collected by the central rings arrive at small
transaxial angles, where parallax is small and DOI is less critical, so the main
requirement is collection efficiency --- where the high density of BGO helps. The
wing rings collect the oblique coincidences at large axial angle, where parallax
is large (so DOI matters) and the photons traverse long paths through the patient
(so scatter, and hence energy resolution, matters) --- the regime where pure CsI
performs well. Placing the dense material where the angles are small and the
high-resolution, DOI-capable material where they are oblique follows the same
principle: the core maximises capture, the wings reconstruct the oblique pairs a
long AFOV adds.

Both materials sit in a single liquid-nitrogen cryostat, and cross-material
coincidences --- one gamma in the BGO core, its partner in a CsI wing --- are
handled naturally, since each crystal independently reports time, energy and a
3D vertex. The main additional complexity is in the readout: BGO and pure CsI
emit at different wavelengths, calling for region-optimised photosensors (or a
broadband choice) and a per-material reconstruction network. This is the most
complex of the configurations considered here.

\medskip
Together with the pure-CsI baseline, these three extensions span a range of
cost and performance --- from an uncooled CsI(Tl) entry-level scanner, through a
density-optimised cryogenic BGO option, to a hybrid that places each material
according to the local incidence angle --- all within the CRYSP architecture and
all aimed at the same low-rate, photon-starved measurement.

\section*{Acknowledgements}
Placeholder.

\bibliographystyle{IEEEtran}
\bibliography{biblio}

\end{document}
